\documentclass[12pt]{article}
\usepackage[canadien]{babel} 
\usepackage[utf8]{inputenc}
\usepackage[T1]{fontenc}
\usepackage{fancyhdr}
\usepackage{graphicx}

\usepackage{enumerate}
\usepackage{amsmath}
\usepackage{amssymb}
\usepackage{subfigure}
\usepackage{amsmath}
\usepackage{amssymb}
\usepackage{bm}
\usepackage{url}
\usepackage[dvipsnames]{xcolor}
\usepackage{todonotes}

% If your paper is accepted, change the options for the package
% aistats2e as follows:
%
%\usepackage[accepted]{aistats2e}
%
% This option will print headings for the title of your paper and
% headings for the authors names, plus a copyright note at the end of
% the first column of the first page.
\setlength{\parindent}{0cm}
\addtolength{\oddsidemargin}{-2cm}
\addtolength{\evensidemargin}{-2cm}
\setlength{\textwidth}{17.78cm}
\addtolength{\topmargin}{-2.25cm}
\setlength{\textheight}{23.24cm}
\addtolength{\parskip}{5mm}
\pagestyle{fancy}

%************
%* COMMANDS *
%************

\input{../math_commands.tex}

\newif\ifexercise
\exercisetrue
%\exercisefalse

\newif\ifsolution
\solutiontrue
%\solutionfalse

\usepackage{booktabs}
\usepackage[ruled,vlined]{algorithm2e}

\usepackage{amsthm}
\theoremstyle{definition}
\newtheorem{exercise}{Question}%[chapter]
\newtheorem{answer}{Answer} % asterisk to remove ordering

\newcommand{\Exercise}[1]{
\ifexercise#1\fi
}

\definecolor{answer}{rgb}{0.00, 0.12, 0.60}
\newcommand{\Answer}[1]{%
  \ifsolution
    {\color{answer}\begin{answer}#1\end{answer}}%
  \fi
}

\newif\ifexercise
\exercisetrue
%\exercisefalse

\newif\ifsolution
\solutiontrue
%\solutionfalse

\usepackage{enumitem}
\newcommand{\staritem}{
\addtocounter{enumi}{1}
\item[$\phantom{x}^{*}$\theenumi]}
\setlist[enumerate,1]{leftmargin=*, label=\arabic*.}

\begin{document}


\fancyhead{}
\fancyfoot{}

\fancyhead[L]{
  \begin{tabular}[b]{l}
    IFT6135-A2026  \\
    Prof: Aishwarya Agrawal \\
  \end{tabular}
}
\fancyhead[R]{
  \begin{tabular}[b]{r}
    Assignment 2 - Theory  \\
    Recurrent Neural Networks and Transformers \\
  \end{tabular}
}
\fancyfoot[C]{- Do not distribute -}

\vspace{1cm}

\shorthandoff{:}
{\textbf{Due Date: November 15th, 2026 at 11:00 pm}}\\


\vspace{-0.5cm}
\underline{Instructions}%
\renewcommand{\labelitemi}{\textbullet}

\renewcommand{\labelitemi}{\textbullet}

\begin{itemize}
\item \emph{This assignment is intensive, start well ahead of time!}
\item \emph{For all questions, show your work.}
\item \emph{Use LaTeX and the template we provide when writing your answers.
You may reuse most of the notation shorthands, equations and/or tables.
See the assignment policy on the course website for more details.}
\item \emph{Submit this assignment on Gradescope. Make sure that you select the pages containing the answer for each question on Gradescope. We will not grade questions without selected pages.}
\item \emph{TAs for this assignment are \textbf{Achille Sowa} and \textbf{Djalil Chikhi}.}
\end{itemize}
\vspace{0.2cm}


\vspace{0.2cm}

% Q1
\begin{exercise}[12] \textbf{Link between Linear Transformer and RNN}
\Exercise{\label{ex:lstm}}

\noindent Let $\vx_1, \dots, \vx_L \in \R^{d}$ be a sequence of $L$ input vectors. A decoder
transformer uses \emph{causal softmax attention}: at time step $t$, the output $\vo_t$ attends
over all the positions $i \le t$ of the sequence,
\begin{equation}
    \vq_t = \mW_Q \vx_t, \qquad \vk_t = \mW_K \vx_t, \qquad \vv_t = \mW_V \vx_t, \qquad
    \vo_t = \sum_{i=1}^{t} \frac{\exp\!\left(\vk_i^\top \vq_t\right)}{\sum_{j=1}^{t} \exp\!\left(\vk_j^\top \vq_t\right)}\, \vv_i,
    \label{eq:softmax-attn}
\end{equation}
where $\mW_Q, \mW_K \in \R^{d \times d}$ and $\mW_V \in \R^{d \times d}$ are learned
projection matrices. Consequently $\vq_t, \vk_t \in \R^{d}$, $\vv_t \in \R^{d}$ and
$\vo_t \in \R^{d}$.

\noindent \emph{Linear attention} replaces the similarity $\exp(\vk_i^\top \vq_t)$ by a dot
product in some feature space, $\phi(\vk_i)^\top \phi(\vq_t)$, for a feature map
$\phi : \R^{d_k} \to \R^{d_\phi}$. In this exercise we consider the simplest case
$\phi = \mathrm{id}$, i.e.\ the similarity is simply $\vk_i^\top \vq_t$.

\begin{enumerate}[label=\textbf{\arabic*.}, leftmargin=*]

\item \textbf{(2 points)} Show that, when $\exp(\vk_i^\top \vq_t)$ is replaced by
$\vk_i^\top \vq_t$ in~\eqref{eq:softmax-attn}, the output can be written as
\begin{equation}
    \vo_t = \frac{\mS_t\, \vq_t}{\vz_t^\top \vq_t},
    \qquad \text{with} \qquad
    \mS_t = \sum_{i=1}^{t} \vv_i \vk_i^\top \in \R^{d \times d},
    \qquad
    \vz_t = \sum_{i=1}^{t} \vk_i \in \R^{d}.
    \label{eq:linear-attn}
\end{equation}

\noindent This can be interpreted as a recurrent neural network (RNN) whose hidden state
$\mS_t$ is updated at every time step, and from which the output is obtained by a linear
projection with $\vq_t$. In practice, the denominator is often dropped, and the output
becomes
\begin{equation}
    \vo_t = \mS_t\, \vq_t.
    \label{eq:linear-rnn}
\end{equation}
This is the convention we follow in the rest of the exercise.

\item \textbf{(3 points)} \emph{Recurrent form.} Using~\eqref{eq:linear-rnn}, write an
algorithm that computes all the outputs $\vo_1, \dots, \vo_L$ from the inputs
$\vx_1, \dots, \vx_L$ in time \emph{linear} in the sequence length $L$.

\emph{Hint:} first derive a recurrence relation expressing $\mS_t$ in terms of $\mS_{t-1}$,
then use it to write your algorithm.
Complete the blanks in the skeleton below.
Complete the blanks in the skeleton below.

\begin{minipage}{\linewidth}
    \begin{algorithm}[H]
    \DontPrintSemicolon
    \KwIn{Inputs $\vx_1, \dots, \vx_L \in \R^{d}$, weights $\mW_Q, \mW_K \in \R^{d \times d}$ and $\mW_V \in \R^{d \times d}$}
    \KwOut{Outputs $\vo_1, \dots, \vo_L \in \R^{d}$}
    Initialization: $\mS_{0} \leftarrow \blank$\;
    \For{$t = 1$ \KwTo $L$}{
      Compute the query, key and value: $\vq_t \leftarrow \mW_Q \vx_t$, $\vk_t \leftarrow \mW_K \vx_t$, $\vv_t \leftarrow \mW_V \vx_t$\;
      %Update the hidden state: $\mS_{t} \leftarrow \blank$\;
      %Compute the output: $\vo_{t} \leftarrow \blank$\;
      Complete this part
    }
    \caption{Linear attention, recurrent form}
    \label{alg:linear-rnn}
    \end{algorithm}
\end{minipage}

\item \textbf{(4 points)} \emph{Parallel form.} On a GPU, it is more efficient to process
many vector products at once as a single matrix product. Stack the inputs and outputs as
rows of matrices,
\[
    \mX = \begin{bmatrix} \vx_1^\top \\ \vdots \\ \vx_L^\top \end{bmatrix} \in \R^{L \times d},
    \qquad
    \mO = \begin{bmatrix} \vo_1^\top \\ \vdots \\ \vo_L^\top \end{bmatrix} \in \R^{L \times d_v},
\]
and let $\mQ = \mX \mW_Q^\top \in \R^{L \times d}$, $\mK = \mX \mW_K^\top \in \R^{L \times d}$
and $\mV = \mX \mW_V^\top \in \R^{L \times d}$. Show that the outputs of~\eqref{eq:linear-rnn}
can be written in matrix form as
\begin{equation}
    \mO = \left( \mQ \mK^\top \odot \mM \right) \mV,
    \label{eq:parallel}
\end{equation}
where $\odot$ denotes the element-wise (Hadamard) product and $\mM \in \{0,1\}^{L \times L}$
is the causal mask, $M_{ti} = 1$ if $i \le t$ and $M_{ti} = 0$ otherwise.

\item \textbf{(3 points)} \emph{Complexity.} Give the time and the average number of sequential steps to compute
all the outputs $\vo_1, \dots, \vo_L$ with each of the three methods below, as a
function of $L$ and $d$ (you may ignore the cost of computing $\mQ$, $\mK$ and $\mV$).
Justify your answers briefly, and comment on the dependence on $L$.

\begin{center}
\begin{tabular}{lcc}
    \toprule
    Method & Time Complexity & Sequential Operations \\
    \midrule
    Softmax attention (Classic transformer) & & \\
    Linear RNN, recurrent form (Question~2) & & \\
    Linear RNN, parallel form~\eqref{eq:parallel} (Question~3) & & \\
    \bottomrule
\end{tabular}
\end{center}

\end{enumerate}




\end{exercise}

\Answer{\hfill

}
% Q2
\begin{exercise}[1-1-4-6-2] (\textbf{Decoding.})
\Exercise{
\label{ex:decoding}
Suppose that we have a vocabulary containing $N$ possible words, including a special token \texttt{<BOS>} to indicate the beginning of a sentence. Recall that in general, a language model with a full context can be written as
\begin{equation*}
    p(w_{1}, w_{2}, \ldots, w_{T} \mid w_{0}) = \prod_{t=1}^{T} p(w_{t} \mid w_{0}, \ldots, w_{t-1}).
\end{equation*}
We will use the notation $\vw_{0:t-1}$ to denote the (partial) sequence $(w_{0}, \ldots, w_{t-1})$. Once we have a fully trained language model, we would like to generate realistic sequences of words from our language model, starting with our special token \texttt{<BOS>}. In particular, we might be interested in generating the most likely sequence $\vw_{1:T}^{\star}$ under this model, defined as
\begin{equation}
    \vw_{1:T}^{\star} = \argmax_{\vw_{1:T}} p(\vw_{1:T} \mid w_{0} = \textrm{\texttt{<BOS>}}).
    \label{eq:joint_prob}
\end{equation}
For clarity we will drop the explicit conditioning on $w_{0}$, assuming from now on that the sequences always start with the \texttt{<BOS>} token.

\begin{enumerate}[label=\arabic{exercise}.\arabic*]
    \item How many possible sequences of length $T+1$ starting with the token \texttt{<BOS>} can be generated in total? Give an exact expression, without the $O$ notation. Note that the length $T+1$ here includes the \texttt{<BOS>} token.
    
    \item To determine the most likely sequence $\vw_{1:T}^{\star}$, one could perform exhaustive enumeration of all possible sequences and select the one with the highest joint probability (as defined in equation \ref{eq:joint_prob}). Comment on the feasibility of this approach. Is it scalable with vocabulary size?

    \item In light of the difficulties associated with exhaustive enumeration, it becomes essential to employ practical search strategies to obtain a reasonable approximation of the most likely sequence.
    
    In order to generate $B$ sequences having high likelihood, one can use a heuristic algorithm called \emph{Beam search decoding}, whose pseudo-code is given in Algorithm~\ref{alg:beam-search-decoding} below

    \begin{minipage}{\linewidth}
    \begin{algorithm}[H]
    \DontPrintSemicolon
    \KwIn{A language model $p(\vw_{1:T} \mid w_{0})$, the beam width $B$}
    \KwOut{$B$ sequences $\vw_{1:T}^{(b)}$ for $b \in \{1, \ldots, B\}$}
    Initialization: $w_{0}^{(b)} \leftarrow \textrm{\texttt{<BOS>}}$ for all $b \in \{1,\ldots,B\}$\;
    Initial log-likelihoods: $l_{0}^{(b)} \leftarrow 0$ for all $b \in \{1,\ldots,B\}$\;
    \For{$t = 1$ \KwTo $T$}{
      \For{$b = 1$ \KwTo $B$}{
        \For{$j = 1$ \KwTo $N$}{
          $s_{b}(j) \leftarrow l_{t-1}^{(b)} + \log p(w_{t} = j \mid \vw_{0:t-1}^{(b)})$\;
        }
      }
      \For{$b = 1$ \KwTo $B$}{
        Find $(b', j)$ such that $s_{b'}(j)$ is the $b$-th largest score\;
        Save the partial sequence $b'$: $\widetilde{\vw}_{0:t-1}^{(b)} \leftarrow \vw_{0:t-1}^{(b')}$\;
        Add the word $j$ to the sequence $b$: $w_{t}^{(b)} \leftarrow j$\;
        Update the log-likelihood: $l_{t}^{(b)} \leftarrow s_{b'}(j)$\;
      }
      Assign the partial sequences: $\vw_{0:t-1}^{(b)} \leftarrow \widetilde{\vw}_{0:t-1}^{(b)}$ for all $b \in \{1, \ldots, B\}$\;
    }
    \caption{Beam search decoding}
    \label{alg:beam-search-decoding}
    \end{algorithm}
    \end{minipage}
    
    What is the time complexity of Algorithm~\ref{alg:beam-search-decoding}? Its space complexity? Write the algorithmic complexities using the $O$ notation, as a function of $T$, $B$, and $N$. Is this a practical decoding algorithm when the size of the vocabulary is large?

    
    \item The different sequences that can be generated with a language model can be represented as a tree, where the nodes correspond to words and edges are labeled with the log-probability $\log p(w_{t}\mid \vw_{0:t-1})$, depending on the path $\vw_{0:t-1}$. In this question, consider the following language model (where the low probability paths have been removed for clarity)

    \hspace*{-2.5em}\begin{minipage}{\linewidth}
    \centering
    \begin{tikzpicture}[scale=0.9,every node/.style={shape=rectangle, inner sep=5pt, align=center}, level 3/.style={sibling distance=5.5em}, level 2/.style={sibling distance=11em}, level 1/.style={sibling distance=22em}, edge from parent path={(\tikzparentnode.south) .. controls +(0,-0.7) and +(0,0.7) .. (\tikzchildnode.north)}, edge from parent/.style={draw, thick, -latex}, logl/.style={draw=none, font=\footnotesize}, label/.style={anchor=west, align=flush left}]
      \node {\texttt{<BOS>}}
        child { node {\texttt{the}}
          child { node {\texttt{green\vphantom{l}}}
            child { node {\texttt{ellipse\vphantom{q}}} edge from parent node[above left,logl] {$-3.6$} }
            child { node {\texttt{rectangle}} edge from parent node[above right,logl] {$-3.3$} } edge from parent node[above left,logl] {$-3.4$} }
          child { node {\texttt{red}}
            child { node {\texttt{circle\vphantom{q}}} edge from parent node[above left,logl] {$-4.8$} }
            child { node {\texttt{rectangle}} edge from parent node[above right,logl] {$-4.5$} } edge from parent node[above right,logl] {$-3.1$} } edge from parent node[above left,logl] {$-1.5$} }
        child { node {\texttt{a}}
          child { node {\texttt{blue}}
            child { node {\texttt{ellipse\vphantom{q}}} edge from parent node[above left,logl] {$-5.4$} }
            child { node {\texttt{square\vphantom{l}}} edge from parent node[above right,logl] {$-5.0$} } edge from parent node[above left,logl] {$-2.5$} }
          child { node {\texttt{red}}
            child { node {\texttt{circle\vphantom{q}}} edge from parent node[above left,logl] {$-4.9$} }
            child { node {\texttt{square\vphantom{l}}} edge from parent node[above right,logl] (logl_lvl3) {$-5.1$} } edge from parent node[above right,logl] (logl_lvl2) {$-2.7$} } edge from parent node[above right,logl] (logl_lvl1) {$-1.8$} };
            
            \node[label, right=1em of logl_lvl3] (logl_lvl3_label) {$\log p(w_{3}\mid \vw_{0:2})$};
            \node[label] (logl_lvl2_label) at (logl_lvl2 -| logl_lvl3_label.west) {$\log p(w_{2}\mid \vw_{0:1})$};
            \node[label] (logl_lvl1_label) at (logl_lvl1 -| logl_lvl3_label.west) {$\log p(w_{1}\mid w_{0})$};
    \end{tikzpicture}
    \end{minipage}
    
    \begin{enumerate}[label=\arabic{exercise}.\arabic{enumi}.\alph*]
        
        \item \emph{Greedy decoding} is a simple algorithm where the next word $\overline{w}_{t}$ is selected by maximizing the conditional probability $p(w_{t}\mid \overline{\vw}_{0:t-1})$ (with $\overline{w}_{0} = \textrm{\texttt{<BOS>}}$)
        \begin{equation*}
            \overline{w}_{t} = \argmax_{w_{t}} \log p(w_{t}\mid \overline{\vw}_{0:t-1}).
        \end{equation*}
        Find $\overline{\vw}_{1:3}$ using greedy decoding on this language model, and its log-likelihood $\log p(\overline{\vw}_{1:3})$.
        
        \item Apply beam search decoding with a beam width $B = 2$ to this language model, and find $\vw_{1:3}^{(1)}$ and $\vw_{1:3}^{(2)}$, together with their respective log-likelihoods.

    \end{enumerate}

    \item Please highlight the primary limitation that stands out to you for each of the discussed methods (greedy decoding and beam search).

\end{enumerate}
}
    
\end{exercise}

\Answer{\hfill

}


% Q3
\begin{exercise}[10-8-12] (\textbf{Temperature in Soft Attention})
\Exercise{\label{ex:temperature-attention}}

Consider the soft-attention mechanism seen in class. Let
$h_1,\ldots,h_{T_x}$ be the encoder states for an input sequence of length
$T_x$. At decoding step $i$, the decoder assigns an alignment score $e_{ij}$
to each encoder state $h_j$. The corresponding attention weight
$\alpha_{ij}$ determines the contribution of $h_j$ to the context vector
$c_i$.

We modify the attention weights by introducing a positive temperature
$\tau_i>0$:
\[
\alpha_{ij}(\tau_i)
=
\frac{\exp(e_{ij}/\tau_i)}
{\sum_{k=1}^{T_x}\exp(e_{ik}/\tau_i)},
\qquad
c_i
=
\sum_{j=1}^{T_x}\alpha_{ij}(\tau_i)h_j.
\]

Here, $i$ denotes the decoding step, $j$ denotes a position in the input
sequence, and $T_x$ is the input sequence length. The attention weights
satisfy $\sum_{j=1}^{T_x}\alpha_{ij}=1$.

In Questions 1 and 2, we consider a fixed decoding step. To simplify the
notation, we write
\[
e_j := e_{ij},
\qquad
\alpha_j := \alpha_{ij},
\qquad
\tau := \tau_i,
\qquad
c := c_i.
\]

\begin{enumerate}

    \item \textbf{Effect of the temperature.}

    \begin{enumerate}[label=(\alph*)]
        \item Consider the alignment scores $e=[2,1,0]$. Compute
        $\alpha(1)$ and determine
        \[
        \lim_{\tau\rightarrow0^+}\alpha(\tau)
        \qquad\text{and}\qquad
        \lim_{\tau\rightarrow\infty}\alpha(\tau).
        \]
        Give the corresponding limiting context vectors and briefly explain
        how $\tau$ affects attention selectivity.

        \item For arbitrary alignment scores, show that
        \[
        \frac{\partial\alpha_j}{\partial e_k}
        =
        \frac{1}{\tau}\alpha_j
        \left(\delta_{jk}-\alpha_k\right),
        \]
        where $\delta_{jk}=1$ if $j=k$ and $0$ otherwise.

        \item Using the previous result, show that
        \[
        \frac{\partial c}{\partial e_k}
        =
        \frac{\alpha_k}{\tau}(h_k-c).
        \]
        Explain what this expression tells us about the effect of increasing
        $e_k$ on the context vector.
    \end{enumerate}

    \item \textbf{Scaling the alignment scores.}

    \begin{enumerate}[label=(\alph*)]
        \item Suppose all alignment scores are multiplied by the same constant
        $s>0$, such that $e'_j=se_j$, while the temperature changes from
        $\tau$ to $\tau'$. Find a relationship between $s$, $\tau$, and
        $\tau'$ such that
        \[
        \alpha'_j(\tau')=\alpha_j(\tau)
        \]
        for every $j$.

        \item If both the scale of the alignment scores and the temperature
        are learned, is observing $\tau=0.2$ sufficient to conclude that the
        model uses highly selective attention? Justify your answer.
    \end{enumerate}

    \item \textbf{Adaptive temperature.}

    We now return to the original notation, where $i$ denotes the decoding step.
Consider two designs:

    Consider two designs:
    \[
    \text{Design A:}\quad \tau_i=\tau,
    \qquad
    \text{Design B:}\quad \tau_i=f(s_{i-1})>0,
    \]
    where Design A uses one learned temperature and Design B predicts the
    temperature from the previous decoder hidden state.

    \begin{enumerate}[label=(\alph*)]
        \item Give one limitation of Design A and one potential advantage of
        Design B.

        \item A colleague claims that a small predicted $\tau_i$ means that
        the model focuses on one encoder state, while a large $\tau_i$ means
        that it combines several encoder states. Using your previous results,
        explain why this interpretation may be incorrect.
    \end{enumerate}

\end{enumerate}
\end{exercise}

\Answer{

}


% Q4
\begin{exercise}[4-2-4] (\textbf{Optimization and Regularization.})
\Exercise{\label{ex:optimization-regularization}}

In this question, we consider two approaches to regularization during
optimization. Let $\theta$ denote the learnable parameters of a neural network
with loss
\[
\mathcal{L}(f(x^{(i)},\theta),y^{(i)}).
\]

One approach is \textit{weight decay}, which modifies the optimizer update
rule. At each iteration, the parameters are updated as
\[
\theta_{t+1}
=
\theta_t-\Delta\theta_t-\lambda\theta_t,
\]
where $\Delta\theta_t$ depends on the optimizer. For SGD, it is the product of
the learning rate and the gradient. For RMSprop and Adam, the update depends on
moving averages. Refer to the lecture slides for the corresponding update
rules.

Another approach is $L_2$ regularization, which modifies the loss function:
\[
\mathcal{L}_{\mathrm{reg}}
=
\mathcal{L}(f(x^{(i)},\theta),y^{(i)})
+
\gamma\|\theta\|_2^2.
\]

\begin{enumerate}

    \item Consider the following statement:
    \begin{quote}
    ``Using $L_2$ regularization with the usual optimizer update is equivalent
    to adding weight decay to the optimizer update.''
    \end{quote}

    Does this statement hold for SGD, RMSprop, and Adam? Provide mathematical
    arguments for your answers.

    \item We now compare $L_1$ and $L_2$ regularization. The $L_1$-regularized
    loss is
    \[
    \mathcal{L}_{\mathrm{reg}}
    =
    \mathcal{L}(f(x^{(i)},\theta),y^{(i)})
    +
    \gamma\|\theta\|_1.
    \]

    Consider a linear regression model with $L_1$ regularization:
    \[
    \hat{\beta}
    =
    \arg\min_{\beta}
    \left\{
    \frac{1}{2}
    \sum_{i=1}^{N}
    \left(
    y_i-\beta_0-\sum_{j=1}^{p}x_{ij}\beta_j
    \right)^2
    +
    \lambda\sum_{j=1}^{p}|\beta_j|
    \right\}.
    \]

    The corresponding model with $L_2$ regularization is
    \[
    \hat{\beta}
    =
    \arg\min_{\beta}
    \left\{
    \frac{1}{2}
    \sum_{i=1}^{N}
    \left(
    y_i-\beta_0-\sum_{j=1}^{p}x_{ij}\beta_j
    \right)^2
    +
    \lambda\sum_{j=1}^{p}|\beta_j|^2
    \right\}.
    \]

    For two parameters $\beta_0$ and $\beta_1$, prove that the unregularized
    squared-error loss has elliptical level sets.

    \item Which of $L_1$ and $L_2$ regularization promotes sparsity? Give an
    informal explanation without going into mathematical details. Which
    regularizer would you choose for feature subset selection?

\end{enumerate}
\end{exercise}

\Answer{

}
% Q5

\begin{exercise}[5-5-5-5]
(\textbf{Paper Review: Vision Transformer (ViT).})


In this question, you are going to write a \textbf{one-page review} of the paper
\href{https://arxiv.org/abs/2010.11929}
{An Image Is Worth 16x16 Words: Transformers for Image Recognition at Scale}.

Your review should have the following four sections:
\textbf{Summary, Strengths, Weaknesses, and Reflections}.
For each section, we provide guiding questions to consider as you read
the paper. Use these questions to structure your review rather than
answering them individually.

\begin{enumerate}[label=(\theexercise.\arabic*)]

    \item \textbf{Summary:}
    \begin{enumerate}
        \item What problem does the paper address?
        \item What are the main contributions?
        \item How does ViT process images using a Transformer,
        and how does this differ from convolutional neural networks?
        \item What are the main experimental results?
        Focus on facts rather than opinions.
    \end{enumerate}

    \item \textbf{Strengths:}
    \begin{enumerate}
        \item What are the main advantages of the proposed approach?
        \item Does the paper introduce a significant conceptual
        or empirical advance?
        \item Are the experiments well designed and convincing?
        \item What makes this approach useful for the research community?
    \end{enumerate}

    \item \textbf{Weaknesses:}
    \begin{enumerate}
        \item What are the main limitations of ViT?
        \item Are there any missing baselines, datasets,
        or experimental comparisons?
        \item Are there design choices that are not sufficiently justified?
        \item How could the proposed method or its evaluation be improved?
    \end{enumerate}

    \item \textbf{Reflections:}
    \begin{enumerate}
        \item How does this paper relate to concepts discussed in class,
        particularly convolutional networks, attention mechanisms,
        and Transformers?
        \item What future research directions could build on this work?
        \item What new ideas or experiments would you be interested
        in exploring?
    \end{enumerate}

\end{enumerate}

This question is subjective, so we will accept a variety of
well-supported answers. You are expected to critically analyze
the paper and offer your own perspective and ideas, beyond what
the paper itself discusses.

\end{exercise}




\end{document}
